\documentclass[11pt]{article}
\usepackage[margin=1in]{geometry}
\usepackage{amsmath,amssymb,amsthm,mathtools,mathrsfs}
\usepackage{graphicx}
\usepackage{microtype}
\usepackage[hidelinks]{hyperref}

\newtheorem{theorem}{Theorem}
\newtheorem{lemma}{Lemma}
\newtheorem{corollary}{Corollary}
\newcommand{\N}{\mathbb N}
\newcommand{\Ix}[1]{\mathcal I_{#1}}
\newcommand{\cA}{\mathscr A}
\newcommand{\cX}{\mathscr X}
\newcommand{\cY}{\mathscr Y}
\newcommand{\cZ}{\mathscr Z}

\title{Regularizing and Coding Summable Ideals\\
\large Candidate full solutions to Questions 6.1 and 6.2 of arXiv:2011.10638}
\author{Autonomous proof candidate for human review}
\date{11 August 2026}

\begin{document}
\maketitle

\begin{center}
\fbox{\parbox{0.91\textwidth}{\textbf{Status.}
Candidate full affirmative solutions, likely valid.  Every summable ideal
generated by an element of $\cZ$ is isomorphic to a single ideal generated by
an element of $\cY$, answering Question 6.1 more strongly than requested.
Moreover, the set $\cA$ is analytic in $\cZ$, answering Question 6.2.}}
\end{center}

\section{Question and answer}

For a positive sequence $x=(x_n)$ with $\sum_nx_n=\infty$, put
\[
 \Ix{x}=\left\{A\subseteq\N:\sum_{n\in A}x_n<\infty\right\}.
\]
Following the notation of \cite{BL}, let
\[
 \cZ=\{x\in c_0:x_n>0\text{ for every }n,\ \sum_nx_n=\infty\},
 \qquad
 \cY=\left\{y\in\cX:\liminf_n\frac{y_{n+1}}{y_n}>0\right\},
\]
and
\[
 \cA=\{x\in\cZ:\Ix{x}\ne\Ix{y}\text{ for every }y\in\cY\}.
\]
Question 6.1 asks whether every $\Ix{x}$ is, up to the ideal-isomorphism
convention used immediately before the question, a finite or countable direct
sum of ideals generated by elements of $\cY$.  Question 6.2 asks whether
$\cA$ is analytic in $\cZ$.

\begin{figure}[ht]
 \centering
 \includegraphics[width=\textwidth]{figures/problem_crop.png}
 \caption{The two concluding questions in the source, PDF page 6.  This
 packet resolves both.}
\end{figure}

For $A\subseteq\N$, write
\[
 D A=\{n-1:n\in A,\ n\ge2\}.
\]
The central observation is the following exact witness formula.

\begin{theorem}\label{thm:main}
For every $x\in\cZ$,
\begin{equation}\label{eq:witness}
 x\in\cA
 \quad\Longleftrightarrow\quad
 \exists A\subseteq\N\quad A\in\Ix{x}\ \text{and}\ DA\notin\Ix{x}.
\end{equation}
Consequently $\cA$ is an analytic subset of $\cZ$.
\end{theorem}

The right side of \eqref{eq:witness} is visibly an existential real witness.
The substance is that no other obstruction to regularizing the weights
exists.

\section{A comparison lemma for summable ideals}

We first isolate the elementary diagonal-selection argument that supplies a
uniform estimate modulo a summable exceptional set.

\begin{lemma}\label{lem:comparison}
Let $u=(u_n)$ and $v=(v_n)$ be positive null sequences.  Then
\[
 \Ix{u}\subseteq\Ix{v}
\]
if and only if there is a constant $C>0$ for which
\begin{equation}\label{eq:exceptional}
 \sum_{\{n:v_n>Cu_n\}}v_n<\infty.
\end{equation}
\end{lemma}

\begin{proof}
If \eqref{eq:exceptional} holds and $A\in\Ix{u}$, split $A$ into the
exceptional set and its complement to obtain
\[
 \sum_{n\in A}v_n
 \leq \sum_{\{n:v_n>Cu_n\}}v_n+C\sum_{n\in A}u_n<\infty.
\]

Conversely, suppose \eqref{eq:exceptional} fails for every $C$.  Inductively
choose disjoint finite sets $F_k$, supported after all previously chosen
indices, such that
\[
 F_k\subseteq\{n:v_n>2^ku_n\},
 \qquad 1\leq\sum_{n\in F_k}v_n\leq2.
\]
This is possible because the relevant $v$-sum is infinite even after a
finite initial segment is removed, and $v_n\to0$ lets the last selected term
be at most $1$.  For $A=\bigcup_kF_k$,
\[
 \sum_{n\in A}u_n
 <\sum_{k\ge1}2^{-k}\sum_{n\in F_k}v_n
 \leq2\sum_{k\ge1}2^{-k}<\infty,
\]
whereas $\sum_{n\in A}v_n\geq\sum_k1=\infty$.  Thus
$A\in\Ix{u}\setminus\Ix{v}$, proving the contrapositive.
\end{proof}

\section{Backward-shift invariance is exactly regularizability}

\begin{lemma}\label{lem:regularize}
For $x\in\cZ$, the following are equivalent:
\begin{enumerate}
 \item\label{it:regular} there exists $y\in\cY$ with $\Ix{x}=\Ix{y}$;
 \item\label{it:shift} $A\in\Ix{x}$ implies $DA\in\Ix{x}$ for every
 $A\subseteq\N$.
\end{enumerate}
\end{lemma}

\begin{proof}
Assume \ref{it:regular}, and choose $\eta>0$ and $N$ such that
$y_n\geq\eta y_{n-1}$ for $n\geq N$.  If $A\in\Ix{y}$, then
\[
 \sum_{n\in A,\ n\geq N}y_{n-1}
 \leq\eta^{-1}\sum_{n\in A,\ n\geq N}y_n<\infty.
\]
The omitted initial part is finite, so $DA\in\Ix{y}$.  Equality of the two
ideals gives \ref{it:shift}.

Now assume \ref{it:shift}.  Define $v_1=x_1$ and $v_n=x_{n-1}$ for $n\ge2$.
Then \ref{it:shift} says precisely that $\Ix{x}\subseteq\Ix{v}$ (the first
coordinate is harmless).  Lemma~\ref{lem:comparison} supplies $C>0$ such
that, for
\[
 S=\{n\ge2:x_{n-1}>Cx_n\},
 \qquad \sum_{n\in S}x_{n-1}<\infty.
\]
Choose $0<\kappa<\min\{1,C^{-1}\}$ and recursively set
\begin{equation}\label{eq:envelope}
 y_1=x_1,
 \qquad y_n=\max\{x_n,\kappa y_{n-1}\}\quad(n\ge2).
\end{equation}
Thus $y_n/y_{n-1}\geq\kappa$ for every $n\ge2$.

It remains to check that the regularization changes the weights by a
summable amount.  Put $d_n=y_n-x_n\geq0$.  If $n\notin S$, then
$x_{n-1}\leq Cx_n$, so \eqref{eq:envelope} gives
\[
 d_n
 =\max\{0,\kappa(x_{n-1}+d_{n-1})-x_n\}
 \leq\kappa d_{n-1},
\]
because $\kappa C<1$.  If $n\in S$, the same formula gives
$d_n\leq\kappa d_{n-1}+\kappa x_{n-1}$.  Hence, in all cases,
\begin{equation}\label{eq:drecurrence}
 d_n\leq\kappa d_{n-1}+\kappa x_{n-1}\mathbf1_S(n).
\end{equation}
Summing \eqref{eq:drecurrence} up to an arbitrary index $M$ and using
$d_1=0$ yields
\[
 (1-\kappa)\sum_{n=1}^M d_n
 \leq\kappa\sum_{n\in S}x_{n-1}.
\]
Therefore $\sum_nd_n<\infty$.  Since $y=x+d$ with $d$ summable,
\[
 \sum_{n\in A}x_n<\infty
 \quad\Longleftrightarrow\quad
 \sum_{n\in A}y_n<\infty
\]
for every $A\subseteq\N$.  Also $y\geq x$ makes $\sum_ny_n=\infty$, and
$y_n\to0$ follows from $x_n\to0$ and $d_n\to0$.  Thus $y\in\cY$ and
$\Ix{x}=\Ix{y}$.
\end{proof}

Lemma~\ref{lem:regularize} immediately proves the equivalence in
\eqref{eq:witness}: membership in $\cA$ is the negation of
item~\ref{it:regular} of Lemma~\ref{lem:regularize}, hence the negation of its
universal shift condition in item~\ref{it:shift}.

\section{Every summable ideal has one regular summand}

We now obtain a stronger affirmative answer to Question 6.1.  As in the
paragraph preceding that question in \cite{BL}, equality with a direct sum is
understood after an isomorphic identification of the underlying countable
sets.

\begin{theorem}\label{thm:isomorphic}
For every $x\in\cZ$, there are a permutation $\sigma$ of $\N$ and a sequence
$y\in\cY$ such that, with $u_n=x_{\sigma(n)}$,
\[
 \Ix{u}=\Ix{y}.
\]
Consequently $\Ix{x}$ is isomorphic to the one-summand direct sum $\Ix{y}$.
\end{theorem}

\begin{proof}
Because $x_n\to0$ and all terms are positive, its terms can be listed in
nonincreasing order: greedily select a largest unused term.  A largest term
exists at every stage, and every original index is eventually selected since
only finitely many terms are at least as large as its positive value.  This
gives a permutation $\sigma$ for which
\[
 u_1\geq u_2\geq\cdots>0,
 \qquad u_n\to0,
 \qquad \sum_nu_n=\infty.
\]

Let
\[
 S=\{n\ge2:u_{n-1}>2u_n\}.
\]
If $n_1<n_2<\cdots$ enumerate $S$, monotonicity gives
\[
 u_{n_{j+1}-1}\leq u_{n_j}<\tfrac12u_{n_j-1}.
\]
Thus the predecessor weights at the large drops form a geometrically
decreasing sequence, and
\[
 \sum_{n\in S}u_{n-1}<\infty.
\]
Apply Lemma~\ref{lem:comparison} to $u_n$ and
$v_1=u_1$, $v_n=u_{n-1}$ for $n\ge2$, with $C=2$.  It follows that
$\Ix{u}\subseteq\Ix{v}$, or equivalently that $\Ix{u}$ is invariant under
$D$.  Lemma~\ref{lem:regularize} now produces $y\in\cY$ with
$\Ix{u}=\Ix{y}$.

Finally, $A\mapsto\sigma[A]$ is an isomorphism from $\Ix{u}$ onto $\Ix{x}$.
Hence one regular summand suffices.
\end{proof}

\section{Descriptive-set verification}

Identify $\mathcal P(\N)$ with Cantor space $2^{\N}$.  For $m\ge1$, define
continuous functions on $2^{\N}\times\cZ$ by
\[
 s_m(A,x)=\sum_{\substack{n\leq m\\n\in A}}x_n,
 \qquad
 t_m(A,x)=\sum_{\substack{2\leq n\leq m\\n\in A}}x_{n-1}.
\]
Then the witness set
\[
 W=\left\{(A,x):
 \sum_{n\in A}x_n<\infty\ \text{and}\
 \sum_{n\in DA}x_n=\infty
 \right\}
\]
is Borel, since it equals
\[
 \left(\bigcup_{K\ge1}\bigcap_{m\ge1}\{s_m\leq K\}\right)
 \cap
 \left(\bigcap_{K\ge1}\bigcup_{m\ge1}\{t_m>K\}\right).
\]
By Theorem~\ref{thm:main},
\[
 \cA=\{x\in\cZ:\exists A\in2^{\N}\ (A,x)\in W\}.
\]
This is the projection of a Borel subset of the product of two Polish spaces,
and is therefore analytic.  This completes the proof.

\section{Checks and scope}

The comparison lemma is sharp enough to handle sparse, arbitrarily deep
drops in the weights: only the total mass of their predecessors must be
summable.  The geometric envelope \eqref{eq:envelope} then propagates each
repair with ratio $\kappa$, while \eqref{eq:drecurrence} prevents the total
repair from diverging.

Theorem~\ref{thm:isomorphic} uses the isomorphism interpretation explicitly
adopted in the source paragraph introducing Question 6.1.  If the displayed
equality there were instead read as literal equality on the original fixed
coordinates, one summand would be exactly the false assertion already
disproved by the source; a direct sum on a different countable base necessarily
requires an identification.

A bounded search of the run indexes, the exact question text, arXiv records
on summable ideals, and papers indexed as citing \cite{BL} located no later
resolution of either concluding question.  Priority is not asserted;
specialist review is required before dissemination.

\begin{thebibliography}{9}
\bibitem{BL}
M.~Balcerzak and P.~Leonetti,
\emph{Convergent subseries of divergent series},
Results Math. \textbf{77} (2022), article 84;
arXiv:2011.10638.
\end{thebibliography}

\end{document}
